% ===================================================================
%  সারণি নং 6 — বাড়ির ভিতর, আসবাব, খাবার, আলো।
%
%  Ceci n'est PAS une transcription : c'est une traduction, faite en
%  2026, en bengali d'aujourd'hui. Voir texto/bn/10-sarani-01.tex
%  pour la consigne, qui vaut pour les seize fichiers.
% ===================================================================

%%K t06-apar-1 apar
\VUsaut{6.0mm}
\VUtitre{15.0pt}{60}{সারণি নং 6}
\VUsaut{1.4mm}
\VUtitre{11.4pt}{0}{বাড়ির ভিতর, আসবাব, খাবার,}
\VUsaut{0.4mm}
\VUtitre{11.4pt}{0}{আলো।}
\VUsaut{2.2mm}

%%K t06-apar-2 apar
বাড়ি তৈরি হয়ে গেছে। জনের কাকা নিজের নতুন বাড়িতে এসে বসেছেন, যার বিন্যাস
আমরা জানিই। এবার দেখা যাক তিনি সেটিকে কীভাবে সাজিয়েছেন। প্রথমে আমরা দেখব:

%%K t06-tit-1 sub
\VUpk{10.2pt}{40}{শোবার ঘর।}
\VUsaut{3.0mm}

%%K t06-c1-01-1 p
1. --- আমরা একটু বেসময়ে এসে পড়েছি, কারণ কাজের মেয়েটি \VUgras{ঘর}
পরিষ্কারে লেগে আছে।

%%K t06-c1-02-1 p
2. --- \VUgras{হাতমুখ ধোয়ার চৌকির} \textsuperscript{(1)} উপরে এদিক-ওদিক সেই
জিনিসগুলি পড়ে আছে যা দিয়ে ধোয়া, চুল আঁচড়ানো ও সাজা হয়। সেগুলি হল
\VUgras{গামলা} \textsuperscript{(2)} ও \VUgras{জলের জগ} \textsuperscript{(3)},
\VUgras{চুলের ব্রাশ} \textsuperscript{(4)}, \VUgras{চিরুনি} \textsuperscript{(5)},
\VUgras{সাবান} \textsuperscript{(6)} ও \VUgras{স্পঞ্জ} \textsuperscript{(7)}; আর
শেষে কুলকুচোর জলের একটি ছোট \VUgras{শিশি} \textsuperscript{(8)}।

%%K t06-c1-03-1 p
3. --- মাটিতে, চৌকির সামনে, এই যে নোংরা জলের \VUgras{বালতি}
\textsuperscript{(9)}।

%%K t06-c1-04-1 p
4. --- \VUgras{বাইরের ঘরে} \textsuperscript{(10)} আমরা কয়েকটি \VUgras{আঁকশি}
\textsuperscript{(13)} দেখতে পাই আর দুটি \VUgras{ছাতা} \textsuperscript{(11)}, যা
\VUgras{ছাতাদানিতে} \textsuperscript{(12)}।

%%K t06-c1-05-1 p
5. --- দরজার অন্য দিকে \VUgras{আয়নাওয়ালা আলমারি} \textsuperscript{(14)}
দাঁড়িয়ে, যাতে \VUgras{কাপড়চোপড়} \textsuperscript{(15)} রাখা থাকে।

%%K t06-c1-06-1 p
6. --- \VUgras{পরিচারিকা} \textsuperscript{(16)} \VUgras{বিছানা}
\textsuperscript{(17)} পাতছে, এই ফাঁকে তার আলাদা আলাদা অংশ দেখে নেওয়া যাক।
\VUgras{খাটের} \textsuperscript{(20)} উপরে একটি ভালো \VUgras{স্প্রিং-গদি}
\textsuperscript{(21)} আছে, যার উপরে জাস্টিন এখনই পশমের \VUgras{গদি}
\textsuperscript{(22)} পাতবে। তারপর সে চেয়ার থেকে দুটি \VUgras{চাদর}
\textsuperscript{(23)} ও \VUgras{কম্বল} \textsuperscript{(24)} তুলে নেবে। তারপর সে
বিছানার শিয়রে \VUgras{লম্বা বালিশ} \textsuperscript{(25)} ও \VUgras{বালিশ}
\textsuperscript{(26)} রাখবে, যা এখন \VUgras{লেপের} \textsuperscript{(27)} সঙ্গে
\VUgras{বিছানার সামনের গালিচার} \textsuperscript{(32)} উপরে পড়ে আছে। সে
\VUgras{পরদায়} \textsuperscript{(19)} ও \VUgras{চন্দ্রাতপে} \textsuperscript{(18)}
পালকের ঝাড়ন বুলিয়ে নেয়, আর অর্ধেক কাজ হয়ে যায়।

%%K t06-c1-07-1 p
7. --- তারপর তাকে \VUgras{বিছানার পাশের টেবিল} \textsuperscript{(28)} একটু
গোছাতে হবে, \VUgras{অ্যালার্ম ঘড়িতে} \textsuperscript{(29)} দম দিতে হবে,
\VUgras{রাতের বাতি} \textsuperscript{(30)} তৈরি করতে হবে, আর \VUgras{মোমবাতি}
\textsuperscript{(31)} সরিয়ে মোমদানি সাফ করতে হবে।

%%K t06-tit-2 sub
\VUsaut{5.0mm}
\VUpk{10.2pt}{40}{সেলাইয়ের ঝুড়ি ও অগ্নিকুণ্ডের সরঞ্জাম।}
\VUsaut{3.0mm}

%%K t06-c1-08-1 p
8. --- \VUgras{সেলাইয়ের ঝুড়ি} \textsuperscript{(34)}, যা \VUgras{টুলের}
\textsuperscript{(33)} পাশে রাখা, তারও গোছানোর বড় দরকার। তাকে \VUgras{সূচিকর্ম}
\textsuperscript{(35)} ভাঁজ করতে হবে, \VUgras{সেলাইয়ের টেবিলে}
\textsuperscript{(38)} \VUgras{আঙুলি}, \VUgras{সুতো} \textsuperscript{(36)},
\VUgras{সুচ} \textsuperscript{(37)} ও \VUgras{কাঁচি} \textsuperscript{(39)} রাখতে
হবে, আর টেবিলের ঢাকনা \VUgras{চাবি} \textsuperscript{(40)} দিয়ে বন্ধ করতে হবে।

%%K t06-c1-09-1 p
9. --- মাঝখানে রাখা \VUgras{গোল টেবিলে} \textsuperscript{(41)} জাস্টিন
\VUgras{সুরাহির} \textsuperscript{(42)} জল বদলাবে। সে \VUgras{চিনিদানিতে}
\textsuperscript{(43)} \VUgras{চিনি} ঢালবে, \VUgras{চিনির চিমটে}
\textsuperscript{(44)} সাফ করবে ও \VUgras{কাপড়ের ব্রাশ} \textsuperscript{(46)}
সরিয়ে দেবে।

%%K t06-c1-10-1 p
10. --- ঘরের ডান দিকটা তাকে কম কাজ দেবে, কারণ \VUgras{আরামশয্যা}
\textsuperscript{(47)} ও তার \VUgras{গদি} \textsuperscript{(48)} নিজের জায়গায়
আছে, আর যেহেতু ঠান্ডা পড়েনি, \VUgras{অগ্নিকুণ্ডের} \textsuperscript{(49)}
সরঞ্জামে কিছু নাড়াচাড়া হয়নি। \VUgras{বেলচা} \textsuperscript{(50)},
\VUgras{চিমটে} \textsuperscript{(51)}, \VUgras{কাঠের ঠেকনা} \textsuperscript{(55)}
ও \VUgras{ছাইয়ের পরদা} \textsuperscript{(56)} পরিষ্কার; \VUgras{আগুনের পরদা}
\textsuperscript{(54)} নাড়ানো হয়নি আর \VUgras{কাঠের বাক্স}
\textsuperscript{(52)} \VUgras{জ্বালানিতে} \textsuperscript{(53)} ভরা।

%%K t06-noto-1 noto
\VUnotes{143.2mm}{%
(*) Babo = ছোট্ট শিশু; ইংরেজি \textit{baby}, ফরাসি \textit{bébé},
ইতালীয় \textit{bambino}।%
}

%%K t06-noto-2 noto
\VUnotes{143.2mm}{%
(*) Lambrequino = ফরাসি \textit{lambrequin}, স্প্যানিশ
\textit{lambrequines}, পর্তুগিজ \textit{lambrequins}, এমনকি জার্মান ও
ইংরেজিতেও \textit{lambrequins} --- অর্থাৎ জার্মান, ইংরেজি, ফরাসি,
পর্তুগিজ ও স্প্যানিশ, সবেতেই একই।%
}

%%K t06-c1-11-1 p
11. --- তাকে এখনও \VUgras{অগ্নিকুণ্ডের ঘড়ি} \textsuperscript{(57)},
\VUgras{মোমদানি} \textsuperscript{(58)} ও \VUgras{সাজানো রেকাবি}
\textsuperscript{(59)} ঝাড়তে হবে, তারপর \VUgras{আয়না} \textsuperscript{(60)}
মুছতে হবে।

%%K t06-c1-12-1 p
12. --- রইল শিশুর \VUgras{দোলনা} \textsuperscript{(61)} (সেই babo (*)), সেটি তো
আগে থেকেই তৈরি।

%%K t06-tit-3 sub
\VUsaut{5.0mm}
\VUpk{10.2pt}{40}{বৈঠকখানা।}
\VUsaut{3.0mm}

%%K t06-c2-01-1 p
1. --- এবার এই যে বৈঠকখানা। এটি খুব সুন্দর \VUgras{ঘর}। ডান কোণের
অগ্নিকুণ্ডটি একটি সূক্ষ্ম \VUgras{মূর্তি} \textsuperscript{(1)} ও দুটি সুন্দর
\VUgras{ফুলদানিতে} \textsuperscript{(3)} সাজানো, আর মখমলের \VUgras{ঝালরে}
\textsuperscript{(2)} (*) ঢাকা।

%%K t06-c2-02-1 p
2. --- অগ্নিকুণ্ডের ফুলকি যেন গালিচায় আগুন না লাগায়, তাই তার সামনে তামার
একটি \VUgras{আগুন-জাফরি} \textsuperscript{(4)} রাখা হয়েছে।

%%K t06-c2-03-1 p
3. --- তার পাশে একটি সুরুচিসম্মত \VUgras{লেখার টেবিল} \textsuperscript{(5)} খোলা
দাঁড়িয়ে। এখানেই \VUgras{গৃহকর্ত্রী} \textsuperscript{(23)} নিজের চিঠি লেখেন।

%%K t06-c2-04-1 p
4. --- জানালায় \VUgras{জালি-পরদা} \textsuperscript{(6)} লাগানো, যা
\VUgras{দড়ি} \textsuperscript{(8)} বেঁধে রাখে, আর সেই দড়ি \VUgras{আঁকশিতে}
\textsuperscript{(7)} আটকানো, আর তার উপরে ভারী কাপড়ের পরদা, যার উপরে একটি
\VUgras{ঝালরপট্টি} \textsuperscript{(9)}।

%%K t06-c2-05-1 p
5. --- জানালার কুলুঙ্গিতে একটি \VUgras{টবচৌকিতে} \textsuperscript{(11)} সবুজ
\VUgras{গাছ} \textsuperscript{(10)}, একটি চমৎকার তাল, যার পাতা প্রায়
\VUgras{তেলের বাতি} \textsuperscript{(12)} পর্যন্ত পৌঁছায়।

%%K t06-c2-06-1 p
6. --- \VUgras{বাতির ঢাকনি} \textsuperscript{(13)} মারিয়া বানিয়েছে, সেই মনোহর
\VUgras{তরুণীই} \textsuperscript{(14)} যে \VUgras{পিয়ানোয়} \textsuperscript{(15)}
বসে আছে। \VUgras{টুলে} \textsuperscript{(16)} সোজা হয়ে বসে সে একটি সুর অভ্যাস
করছে। সে খুব দক্ষ ও খুব পরিপাটি, যেমনটি তার \VUgras{স্বরলিপির আলমারি}
\textsuperscript{(17)} থেকে বোঝা যায়: তাতে পুরো শৃঙ্খলা।

%%K t06-tit-4 sub
\VUsaut{5.0mm}
\VUpk{10.2pt}{40}{দুষ্টু।}
\VUsaut{3.0mm}

%%K t06-c2-07-1 p
7. --- তার ভাই পল একটি \VUgras{বই} \textsuperscript{(19)} খুঁজছে,
\VUgras{বইয়ের আলমারিতে} \textsuperscript{(18)}। সে এমন \VUgras{ছেলে}
\textsuperscript{(20)} যে চঞ্চল ও একেবারে অবাধ্য। খুব উঁচুতে রাখা বইটি পর্যন্ত
পৌঁছাতে আলমারিতে ঝুলে সে তার উপরে রাখা \VUgras{আবক্ষ মূর্তি}
\textsuperscript{(21)} ফেলে দেওয়ার ঝুঁকি নিচ্ছে।

%%K t06-c2-08-1 p
8. --- এই দুষ্টুমি করার জন্য সে এই সুযোগ নিচ্ছে যে তার মা এক সখীর সঙ্গে কথায়
মশগুল, এক \VUgras{মহিলা} \textsuperscript{(22)} যিনি কয়েক দিন তাঁর সঙ্গে
কাটাতে এসেছেন। মা \VUgras{সোফায়} \textsuperscript{(24)} বসে আছেন,
\VUgras{আরামকেদারার} \textsuperscript{(25)} পাশে, আর তাঁর সখী সেই
\VUgras{চেয়ারে} \textsuperscript{(27)}, যার আমরা কেবল \VUgras{পিঠটা}
\textsuperscript{(26)} দেখতে পাই।

%%K t06-c2-09-1 p
9. --- দেয়ালে টাঙানো বড় \VUgras{ফ্রেমে} \textsuperscript{(30)} এক আত্মীয়ার
\VUgras{প্রতিকৃতি} \textsuperscript{(29)}। টেবিলে রাখা \VUgras{অ্যালবামে}
\textsuperscript{(32)} বাকি পরিবারের \VUgras{ছবি} \textsuperscript{(33)} জড়ো
হয়েছে। ছেলেমেয়েরা একটু আগেই সেটি উলটেপালটে দেখছিল, কারণ \VUgras{চেয়ারগুলি}
\textsuperscript{(31)} এখনও টেবিলের চারপাশে জড়ো।

%%K t06-c2-10-1 p
10. --- চিনামাটির \VUgras{চীনা ফুলদানি} \textsuperscript{(34)}, যা
\VUgras{কাগজ-কাটা ছুরির} \textsuperscript{(35)} পাশে রাখা, মারিয়া নিজের
নামদিনে পেয়েছিল, আর ঘরের কোণ ভরে রাখা \VUgras{ভাঁজ করা পরদা}
\textsuperscript{(36)} তার কাকা চীন থেকে এনেছিলেন।

%%K t06-c2-11-1 p
11. --- সুন্দর \VUgras{ছবিতে} \textsuperscript{(37)} ঝলমলে, যা \VUgras{দেয়ালে}
\textsuperscript{(42)} টাঙানো, ছাদের \VUgras{ঝাড়লণ্ঠনে} \textsuperscript{(38)}
আলোকিত, যার \VUgras{বৈদ্যুতিক বাতি} \textsuperscript{(39)} সন্ধ্যায় এত আনন্দিত
জ্বলে, এই বৈঠকখানা বড় মনোরম লাগে। মেঝেয় পাতা কোমল \VUgras{গালিচা}
\textsuperscript{(40)} আর এত কাজের \VUgras{বৈদ্যুতিক ঘণ্টি} \textsuperscript{(41)}
সেটিকে পুরোপুরি আরামদায়ক করে তোলে।

%%K t06-tit-5 sub
\VUsaut{3.6mm}
\VUpk{10.2pt}{40}{খাবারঘর।}
\VUsaut{3.0mm}

%%K t06-c3-01-1 p
1. --- খাবারের ঘণ্টা বেজে গেছে। মারিয়া ছাড়া সারা পরিবার টেবিলের চারপাশে
জড়ো। খাবারঘরটি একটি চমৎকার মাটির \VUgras{উনুনে} \textsuperscript{(1)} সুখকর
গরম, যার \VUgras{নল} \textsuperscript{(2)} সরু।

%%K t06-c3-02-1 p
2. --- এই ঘরে আসবাব বেশি নেই। দেয়ালের গায়ে, \VUgras{টুলের}
\textsuperscript{(4)} উপরে, একটি সাজানো ছোট \VUgras{ফোয়ারা} \textsuperscript{(3)}
ঝুলছে। কাছেই \VUgras{সাজানোর আলমারি} \textsuperscript{(5)}, যার উপরে ঠান্ডা
পদ, মিষ্টান্ন ও সেই পাত্রগুলি রাখা হয় যেগুলির এখন দরকার নেই।

%%K t06-c3-03-1 p
3. --- তাই উপরের তাকে আমাদের চোখে পড়ে একটি \VUgras{ঝলসানো মুরগি}
\textsuperscript{(7)}, একটি \VUgras{মাছ} \textsuperscript{(8)} ও একটি
\VUgras{বাটি} \textsuperscript{(10)} \VUgras{ফলের} \textsuperscript{(9)}। নিচে
\VUgras{পনির} \textsuperscript{(11)}, \VUgras{রুটি} \textsuperscript{(12)} ও
\VUgras{তেল-সিরকার পাত্র} \textsuperscript{(13)}, যাতে সালাদ সাজানোর সব কিছু
থাকে: তেল, সিরকা, \VUgras{নুন} \textsuperscript{(14)} ও \VUgras{গোলমরিচ}
\textsuperscript{(15)}। শেষে সবচেয়ে নিচের তাকে একটি \VUgras{কেক}
\textsuperscript{(17)}, \VUgras{সরষের পাত্রের} \textsuperscript{(16)} পাশে।

%%K t06-c3-04-1 p
4. --- খাবারঘরের ঘড়ি, যাকে \VUgras{দেয়ালঘড়ি} \textsuperscript{(20)} বলে, খুব
অদ্ভুত আকারের। সেটি কাঠের একটি চৌকাঠে বসানো, যাতে একটি \VUgras{বায়ুচাপমাপক}
\textsuperscript{(19)} ও একটি \VUgras{তাপমাপকও} \textsuperscript{(18)} আছে।

%%K t06-tit-6 sub
\VUsaut{4.6mm}
\VUpk{10.2pt}{40}{খাবার পরিবেশন।}
\VUsaut{3.0mm}

%%K t06-c3-05-1 p
5. --- ঘরের চারপাশে দেয়ালে মোম-পালিশ করা ওক কাঠের \VUgras{কাঠের চৌকাঠ}
\textsuperscript{(21)} লাগানো। কাপড়ের একটি \VUgras{দরজার পরদা}
\textsuperscript{(22)} বারান্দার দিকে যাওয়া দরজাটি বেশ ভালো ঢেকে দেয়। খাবারের
পরে পরিবার সেখানেই কফি খেতে যাবে। \VUgras{পেয়ালা} \textsuperscript{(24)} ও
\VUgras{পিরিচ} \textsuperscript{(25)} আগে থেকেই \VUgras{চিনামাটির আলমারিতে}
\textsuperscript{(23)} সাজানো।

%%K t06-c3-06-1 p
6. --- টেবিলের উপরে ছাদ থেকে একটি \VUgras{ঝোলানো বাতি} \textsuperscript{(26)}
বাঁধা; তার খুব জোরালো আলো সন্ধ্যায় সারা খাবারঘরে ছড়িয়ে পড়ে।

%%K t06-c3-07-1 p
7. --- এবার খোদ \VUgras{খাবার টেবিলটি} \textsuperscript{(27)} দেখা যাক। পরিবার
আসার সময় টেবিল পাতা ছিল। \VUgras{টেবিলক্লথের} \textsuperscript{(28)} উপরে
প্রত্যেক খাদকের জায়গায় রাখা ছিল একটি \VUgras{প্লেট} \textsuperscript{(33)},
একটি \VUgras{ন্যাপকিন} \textsuperscript{(29)}, একটি \VUgras{চামচ}
\textsuperscript{(34)}, একটি \VUgras{কাঁটা} \textsuperscript{(35)}, একটি
\VUgras{ছুরি} \textsuperscript{(36)} ও একটি \VUgras{গ্লাস} \textsuperscript{(32)}।
মদের \VUgras{বোতলও} \textsuperscript{(30)} ছিল আর জলের একটি \VUgras{সুরাহি}
\textsuperscript{(31)}।

%%K t06-c3-08-1 p
8. --- \VUgras{পরিচারক} \textsuperscript{(39)} প্রথমে \VUgras{স্যুপের পাত্র}
\textsuperscript{(37)} এনেছিল, যা থেকে এখনও কিছু স্যুপের ভাপ উঠছে। এখন সে প্রথম
\VUgras{পদ} \textsuperscript{(38)} পরিবেশন করছে, মাংসের। তারপর সে সেই পদগুলি
আনবে যেগুলির নাম আমরা করে ফেলেছি; আর শেষে, \VUgras{মিষ্টান্নের} পরে, সে এই
সুযোগ নেবে যে মালিকেরা বারান্দায় চলে গেছেন, আর টেবিল খালি করে খাবারঘর আবার
গুছিয়ে দেবে।

%%K t06-c3-08-2 p
\textit{টীকা।} --- \textit{খাবার সংক্রান্ত রোজকার শব্দ এখানে কেবল মোটামুটি
দেওয়া হয়েছে। এই তালিকা « হোটেল » (নং 13) ও « বাজার » (নং 15) এই দুই
সারণিতে পূর্ণ করা হবে।}

%%K t06-tit-7 sub
\VUsaut{4.6mm}
\VUpk{10.2pt}{40}{রান্নাঘর।}
\VUsaut{3.0mm}

%%K t06-c4-01-1 p
1. --- আধখোলা দরজা দিয়ে যে রান্নাঘরে আমরা উঁকি দিই, সেখানে একটি মনোরম
বিশৃঙ্খলা মেলে। \VUgras{চুল্লির} \textsuperscript{(1)} সামনে, যেখানে
\VUgras{আগুন} \textsuperscript{(3)} চড়চড় করছে, \VUgras{ঝলসানোর যন্ত্র}
\textsuperscript{(5)} বসানো। একটি চমৎকার \VUgras{ঝলসানো মাংস} \textsuperscript{(7)}
\VUgras{শিকে} \textsuperscript{(6)} আছে আর হয়ে আসছে।

%%K t06-c4-02-1 p
2. --- \VUgras{হাপর} \textsuperscript{(8)} ও \VUgras{কয়লার বেলচা}
\textsuperscript{(9)}, যা এইমাত্র কাজে লেগেছে, \VUgras{কাঠের টুকরোসহ}
\textsuperscript{(10)} মাটিতেই ফেলে রাখা হয়েছে।

%%K t06-c4-03-1 p
3. --- অগ্নিকুণ্ডের উপরের তাক গোছানো: \VUgras{লণ্ঠন} \textsuperscript{(11)},
\VUgras{দেশলাইদানি} \textsuperscript{(13)} যাতে \VUgras{দেশলাই}
\textsuperscript{(12)} আছে, \VUgras{মোমদানি} \textsuperscript{(14)} ও
\VUgras{হাতের মোমদানি} \textsuperscript{(15)} তার \VUgras{মোমবাতিসহ}
\textsuperscript{(16)} --- সব নিজের জায়গায়। কিন্তু বড় \VUgras{কয়লার ঝুড়ি}
\textsuperscript{(18)}, যা সেই \VUgras{কয়লায়} \textsuperscript{(17)} ভরা যা
\VUgras{রাঁধুনি} \textsuperscript{(23)} এইমাত্র ভাঁড়ারঘর থেকে এনেছে,
রান্নাঘরের মাঝখানেই পড়ে রয়ে গেছে।

%%K t06-c4-04-1 p
4. --- সত্যি বলতে বেচারিকে তাড়াতাড়ি করতে হবে, যদি দুপুরের খাবার সময়মতো
তৈরি করতে হয়। এখন সে \VUgras{রান্নার উনুনে} \textsuperscript{(19)} একটি
\VUgras{ঝোল} \textsuperscript{(24)} নাড়ছে, যা যেন পুড়ে না যায়।

%%K t06-c4-05-1 p
5. --- \VUgras{তন্দুরে} \textsuperscript{(20)} একটি কেক সেঁকা হচ্ছে, যা সে
ছেলেমেয়েদের বিকেলের জলখাবারের জন্য বানিয়েছে। উনুনের উপরে \VUgras{হাতা}
\textsuperscript{(21)} ও \VUgras{ঝাঁঝরি} \textsuperscript{(22)} ঝুলছে।

%%K t06-tit-8 sub
\VUsaut{4.0mm}
\VUpk{10.2pt}{40}{রান্নার বাসনপত্র।}
\VUsaut{3.0mm}

%%K t06-c4-06-1 p
6. --- ঘরের পিছনে ঝোলানো \VUgras{বাসনের সেট} \textsuperscript{(25)} খুব
পরিষ্কার। তামার সুন্দর \VUgras{ডেকচি} \textsuperscript{(26)}, \VUgras{দুধের জগ}
\textsuperscript{(27)}, \VUgras{ছাঁকনি} \textsuperscript{(28)} ও \VUgras{কুরুনি}
\textsuperscript{(29)} যত্ন করে মাজা হয়েছে।

%%K t06-c4-07-1 p
7. --- \VUgras{নুনদানি} \textsuperscript{(30)} ছোট আলমারির উপরে \VUgras{চায়ের
পাত্র} \textsuperscript{(31)} ও \VUgras{তেলের বড় বোতলের} \textsuperscript{(32)}
পাশে রাখা। \VUgras{দেরাজে} \textsuperscript{(33)} বোধহয় চামচ-কাঁটা ও রান্নার
ছুরি থাকে।

%%K t06-c4-08-1 p
8. --- \VUgras{ঝাঁটা} \textsuperscript{(34)} ও \VUgras{পালকের ঝাড়ন}
\textsuperscript{(35)} এক কোণে। রাঁধুনি এইমাত্র \VUgras{কাটার গুঁড়ি}
\textsuperscript{(36)} কাজে লাগিয়েছে; সেটি সে জানালার পাশে ফেলে রেখেছে।

%%K t06-c4-09-1 p
9. --- একটি \VUgras{হাত মোছার তোয়ালে} \textsuperscript{(37)} দেয়ালে ঝুলছে,
\VUgras{চোঙার} \textsuperscript{(38)} পাশে। রান্নাঘরের এই অংশে \VUgras{ঝাঁঝরি
শিক} \textsuperscript{(39)}, \VUgras{কাটারি} \textsuperscript{(40)} ও
\VUgras{কাটার পিঁড়ি} \textsuperscript{(41)} রাখা হয়। তারপর কাটারির উপরে, একটি
\VUgras{তাকে} \textsuperscript{(42)}, \VUgras{হাঁড়ি} \textsuperscript{(43)},
\VUgras{তরকারির ডেকচি} \textsuperscript{(44)} তার \VUgras{ঢাকনাসহ}
\textsuperscript{(45)}, আর \VUgras{কড়াই} \textsuperscript{(46)}।
\VUgras{ভাজার কড়াই} \textsuperscript{(47)} পাশেই ঝুলছে।

%%K t06-c4-10-1 p
10. --- এই কোণে কেবল তামার \VUgras{কল} \textsuperscript{(48)} ঝলক দেয়।
\VUgras{নর্দমা-চৌকি} \textsuperscript{(49)}, যার উপরে বড় \VUgras{কলসি}
\textsuperscript{(50)} রাখা, নিজের ছোট আলমারিসহ খুব কাজের হবে, যেখানে
\VUgras{সিরকার পিপে} \textsuperscript{(51)} তার \VUgras{টোঁটিসহ}
\textsuperscript{(52)}, \VUgras{আবর্জনার বাক্স} \textsuperscript{(53)} ও
\VUgras{নোংরা বাসন} \textsuperscript{(54)} রাখা হয়।

%%K t06-tit-9 sub
\VUsaut{4.0mm}
\VUpk{10.2pt}{40}{রান্নাঘরে রাখা আরও জিনিস।}
\VUsaut{3.0mm}

%%K t06-c4-11-1 p
11. --- সেই বড় \VUgras{গামলা} \textsuperscript{(55)} যাতে \VUgras{তরকারি}
\textsuperscript{(58)} ধোয়া হয়, আর ঢালাই লোহার \VUgras{কড়াই}
\textsuperscript{(56)}, যা \VUgras{ইটের মেঝের} \textsuperscript{(57)} উপরে রাখা,
বোধহয় সেই আলমারিরই।

%%K t06-c4-12-1 p
12. --- রাঁধুনির উনুনের অভাব নেই, কারণ এই যে টেবিলের কোণে একটি \VUgras{গ্যাসের
উনুন} \textsuperscript{(59)}, যার উপরে একটি ছোট ডেকচিতে জল গরম হচ্ছে। তা থেকে
বেরোনো \VUgras{ভাপ} \textsuperscript{(60)} বলে দেয় যে সেটি ফুটছে নিশ্চয়। এই জল
বোধহয় কফির জন্য, কারণ টেবিলের অন্য প্রান্তে এই যে \VUgras{কফির পাত্র}
\textsuperscript{(64)} ও \VUgras{কফির যাঁতা} \textsuperscript{(63)}।

%%K t06-c4-13-1 p
13. --- গ্যাসের উনুনটি \VUgras{গ্যাসের টোঁটির} \textsuperscript{(62)} সঙ্গে
জোড়া, একটি লম্বা \VUgras{রবারের নলের} \textsuperscript{(61)} মাধ্যমে।

%%K t06-c4-14-1 p
14. --- \VUgras{দিন-মজুর সহায়িকা} \textsuperscript{(66)}, যে রাঁধুনিকে সাহায্য
করতে আসে, রাতের খাবারের তরকারি তৈরি করছে। খোসা ছাড়িয়ে ও ধুয়ে সে সেগুলি
\VUgras{পাত্রে} \textsuperscript{(65)} রেখে দেবে, যতক্ষণ না রান্নার সময় হয়।

%%K t06-tit-10 sub
\VUsaut{4.0mm}
{\centering\fontsize{10.2pt}{10.2pt}\selectfont\textls[40]{\scshape স্নানঘর।} \textit{(নাওয়ার ঘর।)}\par}
\VUsaut{3.0mm}

%%K t06-c5-01-1 p
1. --- বাড়ির প্রধান ঘরগুলি জেনে নিতে এখন একটিই উঁকি বাকি: স্নানঘরের
সরঞ্জাম দেখা। এতে দেরি হবে না, কারণ সেটি বড় নয়, আর আমরা ঝটপট দেখে নেব যে
\VUgras{নাওয়ার টবের} \textsuperscript{(1)} সঙ্গে একটি ভালো \VUgras{জল গরম করার
যন্ত্র} \textsuperscript{(2)} আছে, যে \VUgras{সাবানদানি} \textsuperscript{(3)}
স্নানকারীর হাতের নাগালে, আর যে \VUgras{তোয়ালে-দণ্ড} \textsuperscript{(5)}
দিয়ে \VUgras{তোয়ালে} \textsuperscript{(4)} শুকানো সহজ হয়।

%%K t06-c5-02-1 p
2. --- এই ঘরের দেয়াল \VUgras{ঝকঝকে টালিতে} \textsuperscript{(6)} ঢাকা, যার
ফলে একটি \VUgras{ঝরনা} \textsuperscript{(7)} বসানো সম্ভব হয়েছে, এই ভয় ছাড়াই
যে ব্যবহারের সময় ছিটে সেগুলি নষ্ট করবে। পা ধোয়ার দস্তার একটি ছোট
\VUgras{গামলা} \textsuperscript{(8)} আসবাব পূর্ণ করে।
\VUsaut{6.0mm}
\VUfilet{22mm}
